\section{Step 2 Inspect the caller and unfinished exits}\label{sha-s02}

Open the cumulative starter and locate its marked regions. The supplied caller can invoke individual entries or run a complete memory or file path. Every missing region branches to a routine-specific development exit. The value three means unfinished work; completed kernel entries return zero, while the completed file entry uses statuses zero through two.

\subsection{Observe and compare}
Use the completed entries available at this point; keep their inputs and observation boundary explicit. The surrounding explanation identifies the state being compared.

\subsection{Expected observation}
The default starter returns development status three without claiming a decoded schedule.

\Needspace{9\baselineskip}
\subsection{Check your understanding}
May a development return be treated as a correct digest?

\begin{itemize}
\item \textbf{A} Yes, if output happens to contain zeros.
\item \textbf{B} No; it explicitly marks incomplete work.
\item \textbf{C} Only for an empty file.
\end{itemize}
